\section{Introducción}

Los registros de electroencefalografía (EEG) frecuentemente presentan canales faltantes debido a problemas de contacto, saturación del amplificador, artefactos de movimiento o rechazo por control de calidad. En estudios prácticos, esto no es una inconveniencia menor: los canales faltantes pueden distorsionar interpretaciones espaciales, degradar biomarcadores clínicos como Potenciales Relacionados con Eventos (ERPs) y reducir la confiabilidad de pipelines de decodificación posteriores en Interfaces Cerebro-Computadora (BCI). Por ello, la interpolación de canales ha sido estudiada durante décadas, desde interpolación espacial clásica hasta reconstrucción basada en grafos.

El Procesamiento de Señales de Grafos (GSP) proporciona un marco natural para este problema al representar electrodos como vértices de un grafo y modelar relaciones inter-canal mediante bordes ponderados. La perspectiva general es consistente con la literatura de señales en grafos asociada a Ortega y colaboradores \cite{shuman2013,ortega2018}, donde suavidad, limitación de banda y estructura espectral se utilizan para reemplazar la interpolación puramente geométrica con estimación restringida por grafo.

Aunque se han propuesto muchos métodos en la literatura \cite{narang2013,puy2018,zhang2024bgsrp}, en la práctica sigue siendo difícil compararlos de manera justa. Diferentes trabajos utilizan diferentes máscaras, conjuntos de datos y convenciones de evaluación, y frecuentemente evalúan métodos novedosos contra configuraciones de línea base no optimizadas, inflando afirmaciones de rendimiento. Este estudio busca reducir esa brecha con un benchmark unificado y optimización rigurosa de hiperparámetros.

Los avances recientes en aprendizaje profundo también se han aplicado a la reconstrucción EEG, incluyendo autoencoders convolucionales \cite{banville2021uncovering}, redes neuronales de grafos \cite{defferrard2016} y arquitecturas transformer. Aunque estos enfoques basados en datos pueden lograr rendimiento competitivo, requieren grandes conjuntos de datos etiquetados, son propensos al cambio de dominio entre sujetos y configuraciones de grabación, y carecen de las garantías de convergencia teórica ofrecidas por optimización basada en modelos. En contraste, métodos GSP como TRSS son \emph{sin entrenamiento}: operan directamente en la señal de prueba utilizando solo la topología de electrodos e hipótesis de suavidad físicamente motivadas, haciéndose inmediatamente aplicables a cualquier montaje EEG sin entrenamiento previo. Esta distinción fundamental motiva nuestro enfoque en métodos basados en modelos en este trabajo, reconociendo el aprendizaje profundo como una dirección de investigación complementaria y promisoria.

Las contribuciones principales de este manuscrito son:
\begin{itemize}
\item Un pipeline unificado y reproducible que cubre construcción de grafo, simulación de canales faltantes, reconstrucción y evaluación cuantitativa en múltiples conjuntos de datos (PhysioNet, BCI Competition IV 2a y MNE Sample).
\item Una comparación inter-familia incluyendo líneas base geométricas (Spherical Spline, TPS), análisis de componentes independientes (ICA) e interpolación avanzada basada en grafos con regularización temporal (TRSS, Tikhonov).
\item Optimización extensiva de hiperparámetros usando el framework Optuna para evaluar justa y exhaustivamente todos los métodos---incluyendo líneas base clásicas---bajo condiciones variables (canales faltantes cercanos versus aleatorios).
\item Un análisis fisiológico demostrando la superioridad de la regularización temporal en la preservación de dinámicas de Potenciales Relacionados con Eventos (ERPs) y Densidad Espectral de Potencia (PSD) bajo pérdida severa de canales.
\end{itemize}

Evaluamos exhaustivamente el rendimiento utilizando seis métricas (MAE, RMSE, SNR, DTW, LSD y Coherencia) bajo \ResScenariosNamesEn~condiciones de enmascaramiento dirigidas por escenarios (abarcando proporciones proporcionales y conteos discretos de canales) con ejecuciones repetidas y estadísticas consolidadas de variabilidad.

La contribución principal de este manuscrito es metodológica: establecemos un framework unificado, reproducible y justo para evaluación comparativa de interpolación EEG. En lugar de proponer un algoritmo fundamentalmente novedoso, demostramos que \emph{la optimización rigurosa de hiperparámetros} cierra sustancialmente la brecha entre líneas base clásicas y métodos modernos. Bajo esta evaluación justa, métodos de regularización temporal como TRSS funcionan competitivamente, con fortaleza particular en la preservación de dinámicas fisiológicas (ERPs, contenido espectral) bajo pérdida severa de canales. Este trabajo sirve como punto de referencia para investigación futura, mostrando tanto la promesa como el techo de rendimiento realista de cada clase de método.
